% Uninput working theorem: limits_ab. Parent to assign numbered ID if accepted.
\subsection*{A local resolution retaining the cusp factor}

Let
\[
 X_n=\{fuv=g s^n\}\subset\mathbb C^5,
 \qquad r=g-f,\qquad q=s^n.
\]
There is a proper bimeromorphic local modification
$\widetilde X_n\to X_n$ with smooth total fourfold
and reduced normal-crossings divisor $r=0$.
For fixed $r\ne0$ it resolves the cyclic base change
of a three-branch $q$-normal-crossings chart.
This is a local statement and does not identify its
global smooth fibres with the prescribed tori.

For $n=1$ first blow up $(u,q)$ and then the pullback
of $(f,g)$. The charts of the resulting modification are
\[
\begin{array}{c|c|c}
 \text{coordinates}&\text{map to }X_1&r\\\hline
 (q,U,v,f)&u=qU,\ g=fUv&f(Uv-1)\\
 (u,Q,G,f)&q=uQ,\ v=GQ,\ g=fG&f(G-1)\\
 (u,v,F,g)&q=uvF,\ f=gF&g(1-F).
\end{array}
\]
To derive them, the $q$ chart of the first blow-up
has $u=qU$, hence $g=fUv$; the ideal $(f,g)$ is
already principal there. Its other chart has
$q=uQ$ and equation $fv=gQ$. Blowing up $(f,g)$
gives the last two charts. These are the actual
blow-up charts, with the common factors cancelled
on the closure of the dense torus; no additional
component of the naive total transform is included.
They are smooth. In the first chart $Uv-1$ is a
smooth coordinate along its zero set, and in the
others $G-1$ or $1-F$ is a coordinate. Thus the
displayed products prove the reduced NC assertion
for $r$.

For arbitrary $n$, pull back this proper modification
by $q=s^n$. The first chart remains smooth. The
second is $uQ=s^n$ times the two coordinates $G,f$.
Resolve its $A_{n-1}$ factor by the unit-interval fan.
The third chart is
\[
 uvF=s^n
\]
times the independent coordinate $g$. Its cone lattice
may be written as $N=\mathbb Z^3$ with coordinates
$(a,b,d)$, where the valuations of $u,v,F,s$ are
$a,b,nd-a-b,d$. The cone is the cone over
$n\Delta_2$. Insert all rays $(i,j,1)$ with
$i,j\ge0$ and $i+j\le n$, and triangulate by the
unit triangles
\[
\begin{aligned}
 &(i,j),(i+1,j),(i,j+1) &&(i,j\ge0,\ i+j\le n-1),\\
 &(i+1,j),(i,j+1),(i+1,j+1)
                       &&(i,j\ge0,\ i+j\le n-2).
\end{aligned}
\]
Every determinant is one in absolute value; the
triangles partition the large triangle and give a
proper smooth toric subdivision. Their number is
$n^2$. All new rays have $s$-height one.

The open set $F\ne0$ is the toric open corresponding
to $nd-a-b=0$. Its fan is the unit subdivision of
the edge $i+j=n$, hence the minimal $A_{n-1}$
resolution times $\mathbb C^*$. This is precisely
the second chart's resolution under $Q=vF$.
Where that chart meets the first, the $A$-type
resolution is an isomorphism. These facts glue
the local proper resolutions to a proper
modification of the whole local $X_n$.

The function $F$ is a toric monomial. Above $F=1$
it is submersive: in a regular toric chart a
nonzero monomial has nonzero derivative in a
coordinate with positive exponent. Equivalently,
on $F\ne0$ it is the displayed torus projection.
Thus $r=g(1-F)$ has two smooth transverse
branches in the third chart after resolution.
The second chart has $r=f(G-1)$ independently
of the resolved surface coordinates, and the
first was already checked. This proves smoothness
and reduced normal crossings over $r$ for all $n$.

For $r\ne0$ in the third chart,
$g=r/(1-F)$, so the local fibre is the resolved
threefold $uvF=s^n$ with $F\ne1$.
For $n=1$, its $q$ divisor is $uvF=0$, a literal
three-branch normal-crossings chart. For $n=12$
the standard triangulation gives 144 unimodular
charts above this one triple point. This is the
same local cyclic-cover count used earlier for
each of the two actual cusp triple points; it
is not a new global Euler count. The other
potential singularity, at $f=-r$, is included
in this construction: the first blow-up small-resolves it
for $n=1$, and the second is an isomorphism there
because $f$ is nonzero.

The equation arises from the actual repaired
neighbourhood of B.4.21. There $h=uv-q$ in the
old node chart and $H=h/q$ is holomorphic in the
repaired collar. The unit inverse-normal plumbing
$Hf=r$ is exactly
\[
 f(uv-q)=rq,
 \qquad fuv=(f+r)q.
\]
Putting $g=f+r$ gives $X_n$. Thus the retained
$q$ factor has been derived from the actual
meromorphic function regularized by the
modification, rather than inferred solely from
an abstract normal-line isomorphism.
More precisely, the actual repaired collar, where $H$ is small,
identifies with the third chart near $F=1$ by
\[
 F=(1+H)^{-1},\qquad g=f(1+H),\qquad r=fH.
\]
There $uvF=s^n$ is exactly the resolved equation
$uv=s^n(1+H)$ from B.4.21. The charts near $F=0$
are additional local completion charts; their attachment
to a complementary global Kato piece has not been established.

This local equation and its resolution do not
supply the complementary Kato charts, their
second-period identification, the proper closure
of the surviving exceptional-curve families, or
an origin cap for B.4.19. In particular the
fixed-base obstruction in B.4.20 is not removed
by recognizing a local Mumford triple chart.
The exact chart identities, determinants, edge
incidences and fan counts are saved in
\path{code/ab_push_corner_resolution.py} and its
output.

The full period action is a separate condition. The established
first collar action is $H\mapsto aH$ and $f\mapsto a^{-1}f$;
the product $r=fH$ is fixed. The product $G(q,H)$ from B.4.21
conjugates this on a small collar, but this does not extend that
conjugacy past its zeros $H=-a^{-j}$. If an actual marked period
were to scale the test coordinate $W=1+H=1+r/f$ by $b$, its
required action on $f$ would be
\[
 f\longmapsto\frac{rf}{(b-1)f+br},
\]
not scalar multiplication. The test coordinate has not been
identified with a marked period coordinate, so this identity
records a compatibility requirement, not an obstruction.
